% ============================================================
% DRAFT NOTE (introduction)
% Contains: ML-framed introduction. Messages in order: compression's harm is
% real and uneven (prevalence paragraph); the gap is temporal (after-the-fact
% evaluation); the forecastability question; why it is hard (attribution ->
% paired damage definition); HERALD and the mechanism hypothesis; earliness
% as characterization; contributions. No evaluation-discipline paragraph by
% decision (2026-09-02); that material lives in the setup section.
% Written now: Everything except the results preview.
% Pending: Results preview; F1 reference; prevalence numbers are catastrophe
% rates carried from the prior draft on this same corpus (verify against the
% corpus audit) and should be complemented or replaced by paired-damage
% prevalence once the label audit lands; GSM8K reference-wrong fraction.
% Sources: docs/goal.md; docs/implementation/quality_risk_protocol.md;
% prior HERALD draft (voice, prevalence table)
% ============================================================
\section{Introduction}

KV-cache compression has become a load-bearing optimization for long-context
language-model inference. Memory and decode latency both scale with cache
size, and a growing family of compressors trades cache fidelity for
throughput. At moderate removal ratios these methods largely preserve task
accuracy; at heavier ratios they do not, and the resulting failures are
user-facing and hard to roll back once bad text has been emitted.

The harm is real, and it is uneven. In the corpus we study, generations of
Qwen2.5-7B-Instruct under six compressors and seven ratios, the share of
GSM8K generations that collapse into looping or fail to terminate ranges from
\todo{4.8\%} under ExpectedAttention to \todo{71\%} under KNorm on the same
prompts and ratios, against a \todo{1.2\%} floor with no compression, and it
grows roughly linearly with the ratio, adding about \todo{ten} points per
0.125 step \todo{verify all four against the corpus audit; add paired-damage
prevalence from the label audit}. Averaged over a benchmark these are tail
events; at heavy ratios under a fragile compressor they are the modal
outcome. A serving system cannot trust a single per-compressor accuracy
number, because the variance sits exactly where it hurts.

Today that harm is measured after the fact. Benchmarks average task accuracy
over completed outputs, which smooths over both which generations were
harmed and when the harm began. Output-level monitors fire only once bad text
is visible. The question this paper asks is whether the damage can instead be
forecast while it is happening: during a generation decoding under
compression, from nothing more than the model's own next-token distributions,
can we estimate whether this generation is going to come out worse than it
would have with the full cache?

Two things make the question harder than it looks. The first is attribution.
Language models are often wrong with a full cache; on GSM8K the model we
study produces a wrong final answer on roughly a fifth of prompts with no
compression at all \todo{verify fraction against the corpus label audit}. A
wrong answer is therefore not evidence of compression damage. Damage has to
be defined against the paired counterfactual, the same prompt decoded by the
same model with an uncompressed cache, and this paper does so: a generation
is damaged when its final task quality is strictly lower than its
uncompressed counterpart's, so wrong-anyway answers and compression-induced
lift do not count. The second is the granularity mismatch. The label arrives
once per generation, at the end; the forecast is wanted at every token, from
the start. A detector that is confident only at the final token is a grader,
not a forecast.

HERALD (Hazard Estimation via Real-time Analysis of Logit Distributions) is
built on one hypothesis: that compression harm is visible in the model's own
next-token distributions, in their uncertainty, their instability from step
to step, and the churn among their leading alternatives, before it is visible
in the text. If that holds, a cheap causal model over those statistics,
together with the compression action and the position, can estimate at every
decoding step the probability that the generation ends damaged. The target is
risk rather than magnitude: a probability is what a downstream decision needs
to weigh against the cost of backing off, and a binary outcome is a far
kinder learning problem than signed magnitude on labels that arrive once per
generation. The claim is a current-state one. Compression is already active
and stays active; forecasting what would happen if it were switched on now is
a separate, prospective problem that we do not address.
Figure~\todo{F1 reference} illustrates the setting.

Because the forecast is per token and the label per generation, we
characterize the forecaster by its earliness: at fixed checkpoints into the
generation, how well the risk estimate separates generations that will end
damaged from those that will not, and how well calibrated it is. We report
that curve as the operational description of the method rather than
promising a detection horizon in advance; where the estimate becomes useful,
and where it degrades by task, compressor, and ratio, is part of the result.

\todo{One or two sentences previewing the confirmation result: whether the
hypothesis survives the preregistered gates, at which checkpoint the risk
estimate becomes useful, and where it degrades.}

This paper makes four contributions:

\begin{enumerate}
  \item \texttt{herald-logits}, a public corpus of \todo{32{,}852} compressed
  and uncompressed generations over \todo{764} prompts from four tasks, six
  compressors, and seven ratios, with per-token logit statistics,
  matched-prefix divergence oracles against the paired uncompressed run, and
  final task-quality and failure annotations. Any compression-damage target,
  including ours, can be defined on it without rerunning a model.
  \item A current-state, compression-attributable operational definition of
  quality damage, with a per-step risk target that is well-defined
  throughout a compressed generation and identically zero without
  compression.
  \item A test of the hypothesis that compression damage is forecastable
  online from zero-cost logit statistics, against a comparator that knows
  the compression action and position but not the stream.
  \item An earliness analysis of the resulting forecaster: how many tokens
  into a generation it becomes reliable, overall and for catastrophic versus
  graceful damage, and where it degrades by task, compressor, and ratio.
\end{enumerate}
